\subsection{章末確認問題}
\begin{enumerate}
  \item 「アーキテクチャ」という語の三つの意味を説明せよ。
  \item 利害関係者と関心事の違いを、自分の言葉で説明せよ。
  \item なぜアーキテクチャ記述では複数のビューを使うのか。
  \item ビューとビューポイントの違いを説明せよ。
  \item アーキテクチャ様式、アーキテクチャパターン、参照アーキテクチャを区別せよ。
  \item アーキテクチャ上重要な要求 ASR の例を三つ挙げよ。
  \item アーキテクチャ根拠を記録する理由を説明せよ。
  \item アーキテクチャ評価で、トレードオフを確認することが重要な理由を説明せよ。
  \item アーキテクチャ測度の例を三つ挙げ、それぞれ何を知るためのものか説明せよ。
\end{enumerate}

\begin{statementbox}{解答の方向性}
1 は三義、2 は主体と論点、3 と 4 は関心事別の表現、5 は型・解決法・領域基準の違いで説明する。6 以降は、性能、可用性、セキュリティ、変更容易性、法規制、外部連携を例にできる。
\end{statementbox}
